\documentclass[11pt]{article}
\usepackage[margin=1.1in]{geometry}
\usepackage{fontspec}
\setmonofont{Noto Sans Mono}
\usepackage{amsmath,amssymb,amsthm}
\usepackage{booktabs}
\usepackage{enumitem}
\usepackage{longtable}
\usepackage{array}
\usepackage{calc}
\usepackage{stmaryrd}
\usepackage{xcolor}
\usepackage{graphicx}
\usepackage{multicol}
\usepackage{tikz}
\usetikzlibrary{arrows.meta,positioning,fit,backgrounds}
\usepackage{pgfplots}
\pgfplotsset{compat=1.18}
\usepackage[colorlinks=true,linkcolor=blue!50!black,citecolor=blue!50!black,urlcolor=blue!50!black]{hyperref}
\usepackage{natbib}
\usepackage{listings}
\lstset{basicstyle=\ttfamily\small,columns=fullflexible,
        keepspaces=true,frame=single,framesep=6pt,
        captionpos=b,breaklines=true,xleftmargin=0pt}

\definecolor{okgreen}{HTML}{2F7D32}
\definecolor{openamber}{HTML}{B5651D}
\definecolor{foldgray}{HTML}{4A4A4A}

\newtheorem{theorem}{Theorem}[section]
\newtheorem{lemma}[theorem]{Lemma}
\newtheorem{proposition}[theorem]{Proposition}
\newtheorem{corollary}[theorem]{Corollary}
\newtheorem{conjecture}[theorem]{Conjecture}
\theoremstyle{definition}
\newtheorem{definition}[theorem]{Definition}
\newtheorem{example}[theorem]{Example}
\newtheorem{openproblem}{Open Problem}
\theoremstyle{remark}
\newtheorem{remark}[theorem]{Remark}
\newcommand{\status}[1]{\textcolor{blue!45!black}{\textsf{[\,#1\,]}}}

\newcommand{\NN}{\mathbb{N}}
\newcommand{\CLS}{\mathrm{CL}(S)}
\newcommand{\LF}{L_{F}}
\newcommand{\LO}{\mathrm{lo}}
\newcommand{\bhs}{\mathrm{BeT}}
\newcommand{\W}{\mathrm{W}}
\newcommand{\hnf}{\mathrm{hnf}}
\DeclareMathOperator{\spine}{spine}
\DeclareMathOperator{\arity}{arity}
\DeclareMathOperator{\Drop}{Drop}


\providecommand{\tightlist}{\setlength{\itemsep}{0pt}\setlength{\parskip}{0pt}}
\setlength{\emergencystretch}{3em}
% Alternative titles, if "High-Level Native" is not earning its place:
%   SKIjack: A Language Whose Staging Line Is a Theorem
%   SKIjack: Metacircular Interpretation as a Language Feature over $SKI$
%   SKIjack: Writing Interpreters for a Calculus That Cannot Quote
\title{SKIjack:\\A High-Level Native Language for the $SKI$ Calculus}
\author{N.\ E.\ Davis}
\date{\today}

\begin{document}
\maketitle

\begin{abstract}
Over a homoiconic calculus the designer chooses where quotation happens.
Over a calculus without self-quotation the choice is not available: no
combinator can take another apart, so reification is metalevel by
theorem, and a language's staging boundary is fixed before its syntax
is.  We present SKIjack, a language for writing metacircular
interpreters over the $SKI$ calculus, whose feature set is read off an
existing artifact rather than designed: the 618-atom self-interpreter of
\citet{Davis2027}, decomposed into its fuel loop, its dispatch, its
spine walker and its atoms, with each part promoted to a surface
construct.  Declaring an object language generates that furniture, so
authoring the semantics of a new symbol --- the operation the first
paper performed by hand --- becomes a two-line edit.  Quotation is
available at compile time only, and the checker refuses statically the
positions at which a function would be asked to behave as data.  A scry
namespace is a language feature, with paths as a declared type and
blocking sound under an append-only namespace.  The abstraction is free,
and we show it by equality of terms: six programs compile
byte-identically to the hand-built artifacts, including one in which
nothing is generated, and reproduce their contraction counts.
\end{abstract}

\section{Introduction}\label{sec:intro}

Over a homoiconic calculus (Lisp, Nock) the metacircular tradition gets 
quotation for free and the language designer chooses where evaluation 
stages live. Over a non-homoiconic calculus that choice is taken away: 
Proposition 3.1 says no term can quote a live term, so the compile-time/
runtime line is fixed before the language is designed, and the two 
operations the calculus cannot do for itself, quote and recognize, go to 
the compiler and the virtual machine respectively.  \citet{Davis2027} 
presented a metacircular evaluator for $SKI$ which is written and 
evaluated in the same combinator system.  This capability immediately 
suggests the construction of a battery of computational conveniences and 
abstractions within the same system; to wit, a language.  We present 
SKIjack (pronounced ``skyjack''), a closed-supercombinator macro-like 
language built in and  targeting the non-homoiconic $SKI$ calculus as 
its  instruction set  architecture.  Every named thing compiles to a 
closed term and every name is resolved before anything runs, so the 
language has no runtime environment to address: bracket abstraction has 
already closed what a subject-oriented calculus would have to carry.  
SKIjack is evaluable on any $SKI$ interpreter as a virtual machine.

The asymmetry that forces the staging line is older than either
calculus's programming languages, and naming it early makes the rest of
the design read as consequence rather than taste.  \citet{Bohm1968}
proved that two distinct $\beta\eta$-normal forms can always be told
apart from the outside: some context applies them to arguments and
recovers any two results we please.\footnote{Stated for the
$\lambda$-calculus and for $\beta\eta$-distinct normal forms; see also
\citet[\S10.4]{Barendregt1984}.  The probes this paper uses run on
Scott-encoded data under weak reduction, which is not that setting, so
we cite separability for the asymmetry it names and not as a
completeness result for the method.}  \citet{JayGivenWilson2011} proved
that no combinator can take another apart from the inside.  Both papers
of this pair live in the gap between those two results.  Reading a
compiled term by applying it to fresh markers and seeing which one comes
back is the outside; Proposition 3.1 is the inside; and a VM jet is
permitted to factorise a live term precisely because it is the Turing
machine with the term on its tape, which is to say outside as well.
Quotation goes to the compiler and recognition to the VM not
because that is convenient but because those are the only places the
operations exist.

The \texttt{whnfF} metacircular interpreter and \texttt{UQ} unquoter 
function presented in \citet{Davis2027} are the starting point for our 
work.  That paper discussed extensions to the $SKI$ calculus built on
higher layers of interpreters, introducing an \texttt{Err} error type,
a \texttt{Maybe} partiality type, and discussed implications such as
namespace value retrieval and other innovations.  By authoring a new
symbol at one site, the \texttt{whnfF} approach enables extension and 
modification of the combinator's semantics, permitting it to behave more
like a language than a mere combinator system.  This paper presents
SKIjack and three novel contributions for $SKI$ combinator research:
a two-level execution model with a kernel/userspace split; and a
namespace as a language feature:  a resolver is a core parameter, paths 
are a declared type compared structurally, and the namespace literal 
compiles to a resolver.

\paragraph{Contributions.}
\begin{itemize}[leftmargin=1.4em,itemsep=2pt]
\item A language whose staging boundary is forced rather than chosen,
  and whose checker enforces it before the expander runs: the positions
  that require data --- equality, quotation, a namespace key, a path ---
  refuse functions, so no accepted program can ask for an operation the
  calculus cannot perform (\S\ref{sec:boundary}).
\item Object-language declarations that generate an interpreter's
  furniture --- constructors, case, spine walker, rebuilder, default
  arms and fuel loop --- which makes authoring a semantics at one site a
  two-line edit rather than a term construction
  (\S\ref{sec:interpreters}).
\item A scry namespace as a language feature: typed paths, a resolver as
  a core parameter, a namespace literal that compiles to a resolver, and
  blocking whose soundness rests on the namespace being append-only,
  with the non-neutrality of fuel across a resume stated rather than
  hidden (\S\ref{sec:namespace}).
\item An expander that is a fold over the syntax tree, with bijective
  laws --- $\mathit{expand} \circ \mathit{render} = \mathit{lower}$ and
  $\mathit{lower} \circ \mathit{lift} = \mathit{id}$ --- and a published
  dictionary keyed by structural hash, which is the same key a VM
  matches on (\S\ref{sec:impl}).
\item Evidence that the abstraction costs nothing, by equality of terms:
  six programs compile byte-identically to artifacts built by hand
  before this language existed, and reproduce their contraction counts
  (\S\ref{sec:impl}).
\end{itemize}

\section{Background}\label{sec:background}

\status{One page, cited and not restated. $SKI$ and bracket abstraction;
Scott encoding and the continuation orders that are the ABI; the
artifact of \citet{Davis2027} --- \texttt{whnfF}, the unquoter
\texttt{UQ}, the fuel semantics in which fuel $k$ permits $k$
step-attempts, and Proposition 3.1 --- and the counts this paper
reproduces rather than re-derives. Fix the measure convention here as
well, since four rulers get used on the same object and a reader who
confuses them will think the paper contradicts itself: \texttt{whnfF}
is 618 \emph{atoms}, 1{,}235 tree \emph{nodes}, 348 distinct nodes once
sharing is counted, and 2{,}109 \emph{characters} when printed. Atoms
are the unit for every claim in this paper; the others appear only where
they are named. Everything a reader needs for \S\ref{sec:interpreter}
and nothing more. The separation of credit belongs
here: the first paper establishes that a self-interpreter exists and
what it costs; this one asks what a language owes it.}

\section{Lineage}\label{sec:lineage}

SKIjack is a language entirely based on the $SKI$ combinator calculus,
and in particular on improving the expressiveness of $SKI$ reasoning by
the introduction of a standard machinery in the form of generated
supercombinators.  We survey the historical and technical lineage of
pure combinator computation that informs SKIjack's design.  SKIjack is 
far from the first language or machine concept to lean into
pure combinator computing.  The present paper leans on five classic 
results from Corrado B\"ohm's work:
\begin{enumerate}[nosep]
  \item  \emph{CUCH}.  \citet{BohmGross1966} introduced the CUCH (CUrry-CHurch), a surface for writing programs whose meaning is given by combinators.  The $SKI$ calculus (including SKIjack) operates under a more severe constraint:  while CUCH's designers were able to place its affordances such as quotation wherever they pleased, pure $SKI$ calculus cannot quote expressions (Proposition~3.1 of \citet{Davis2027}).
  \item  \emph{Separability}.  TODO
  \item  \emph{Term-algebra schema}.  A type declaration generates its constructors and its case form by the 
  schema \citet{BohmBerarducci1985} gave for term algebras, in the Scott 
  variant \citep{Scott1962,Jansen2013} rather than the Church one:  an 
  interpreter's step function wants case analysis in place of a fold which 
  fixes the declaration order as the continuation order.
  \item  \emph{First self-compiling compiler}.  The first compiler written in its own language is recognized as \citet{Bohm1954}, his 1951 dissertation at ETH Z\"urich \citep{KnuthPardo1980}.  In the case of SKIjack, the compiler carries out quotation at compile time.
  \item  \emph{Self-interpreter with a normal form}.  Finally, \citet{BerarducciBohm1993} gave a self-interpreter of the $\lambda$-calculus having a normal form.  
\end{enumerate}

Bracket abstraction permits the procedural compilation of a high-level 
applicative language to combinators \citep{Turner1979}; there is a 
related abstraction algorithm \citep{Turner1979b}.  \citet{Hughes1982} 
introduced the notion of a supercombinator, the closed abstraction 
compiled as one unit.  This program effloresced in the ``$SKI$ Machine''
(SKIM), a graph-reduction machine built on the ``Turner set'' of 
combinators, namely, $S$, $K$, $I$, $B$, $C$, $Y$, and $P$
\citep{Johnsson1984,ClarkeEtAl1980} and its successor
\citep{StoyeClarkeNorman1984,Stoye1985}.  (A standard account of the 
whole pipeline has been given by \citet{PeytonJones1987}.)  The SKIM 
tradition operated with a fixed set of recognized combinators executed 
natively in microcode, an accelerated process comparable to the ``jet''
approach we will adopt in this work.  Much of this work was concerned
with the cost of any particular abstraction, in combinators per source 
symbol \citep{Noshita1985,JoyRaywardSmithBurton1985}, and we report atom
counts in line with that tradition.

In the preceding work, the interpreter lived either in a machine (Turner)
or in microcode (SKIM), with the reduction engine as an external given.
The object introduced by \citet{Davis2027} is a 618-atom 
self-interpreter which exists as a term within the language it 
interprets.  Previous supercombinator work has had no reason to write
such an artifact; its self-application question was the compiler
bootstrapping in its own source language, not the reducer reducing
itself, and so the tower, the per-level cost, and the encoding boundary
never arose there.

However, pursuant to the introduction of a self-interpreter as a term 
within the language, new considerations arise.  Staging becomes a design 
problem: Turner's SASL could not quote either, but because no SASL 
program was trying to run another SASL program under a fuel budget, 
there was no obligation to address the lack.  Declaring an object 
language becomes a language feature, since an interpreter's  
constructors, case form, spine walker, rebuilder and fuel loop can be 
generated from the declaration rather than written.  Finally, the native 
compiler owes its output a strict discipline:  expansion is bijective, 
terms lift back to named source, and the dictionary is published and 
keyed by structural hash.  This discipline becomes crucial because the 
emitted term is simultaneously the object of study and the key a VM 
matches on, whereas nobody expected to read the output of a SASL  
compiler.

This last merits an aside.  Our VM is capable of identifying 
supercombinators to arbitrary size and reducing them instantly against
a table of known outcomes in native code (what we call a ``jet'', short
for ``jet-accelerated code'').  Mechanically, such a procedure
corresponds to Turner's extended combinator set and SKIM's microcode.  
However, our theorem does not licence arbitrary native code replacement;
the interpreter jet performs a structural recognition that no
term can perform for itself and thus is not an optimization of an
operation the language has.  The operational consequence of this
prohibits a jet from quoting on behalf of a running term.  The
VM is addressed briefly in \S\ref{sec:runtime}.

\section{The interpreter}\label{sec:interpreter}

The \texttt{whnfF} function forms the centerpiece of the interpreter, 
responsible for reducing terms to weak head normal form.  It
orchestrates the evaluation strategy for the language, ensuring that 
terms are reduced correctly according to the operational semantics.
It also suggests the language features that are necessary to support 
such an evaluation strategy, which will emerge in SKIjack as primitives.
There is also an unquoter function \texttt{UQ}, responsible for 
converting weak head normal forms back into their original syntactic 
representation.  This is necessary for maintaining the correspondence 
between the evaluated terms and their original syntactic forms, ensuring 
that the interpreter's output can be related back to the source code.
In \citet{Davis2027}, we simply named the features by their operational
significance as a shorthand; here they become language elements in their
own right.

Table~\ref{tab:atoms} gives some typed forms with their defining 
equations and with closed $\{S,K,I\}$ terms emitted by the expander
(i.e., compiled output).T hree of the six reduce to
the same term, and two more are a second same term.  This results because
a Scott datum is its own case analysis, so the first constructor of any 
two-constructor type discards the remaining continuation and is $K$, 
whatever the type was called, and a one-field second constructor applies 
one continuation to one field and is the same eight atoms whether the 
field is a numeral or a term.  At the level of terms the type 
distinctions are erased.  SKIjack maintains the type discipline at
compile time (\S\ref{sec:boundary}).  A checker may reject a
program that offers a numeral where a list is wanted, but nothing
downstream of the expander could tell the two apart, because downstream
of the expander they are one object.  (It is also why the dictionary that
names these expansions is keyed by the hash of a term and not by a name.)

\begin{table}[htb]
\centering
\caption{Some significant forms, their definitions, and the terms they 
compile to.  The first constructor of each pair is the same single atom 
and the one-field second constructors are the same eight; only
\texttt{cons}, which carries two fields, is distinct.}\label{tab:atoms}
\footnotesize
\begin{tabular}{@{}l l >{\raggedright\arraybackslash\ttfamily}p{3.35in} r@{}}
\toprule
\textrm{form} & \textrm{definition} & \textrm{expansion} & \textrm{atoms} \\
\midrule
\texttt{nil}            & $\lambda n.\lambda c.\,n$       & K & 1 \\
\texttt{cons}\;$h\;t$   & $\lambda n.\lambda c.\,c\,h\,t$ & S (K (S (K K))) (S (S (K S) (S (K K) (S (K S) (S (K (S I)) K)))) (K K)) & 22 \\
\texttt{nothing}        & $\lambda n.\lambda j.\,n$       & K & 1 \\
\texttt{just}\;$x$      & $\lambda n.\lambda j.\,j\,x$    & S (K K) (S (K (S I)) K) & 8 \\
\texttt{zero}           & $\lambda z.\lambda s.\,z$       & K & 1 \\
\texttt{suc}\;$n$       & $\lambda z.\lambda s.\,s\,n$    & S (K K) (S (K (S I)) K) & 8 \\
\bottomrule
\end{tabular}
\end{table}

The metacircular evaluator uses these object terms to represent the 
basic constructs of the language within itself.  Next, \texttt{whnfF}
itself is composed of 618 atoms built from a few pieces.
Table~\ref{tab:whnfF} inventories \texttt{whnfF}'s parts, smallest 
first, with the term each part compiles to.  By reading the lower block 
from the bottom upwards, the interpreter assembles itself out of the 
upper one.
\begin{multicols}{2}
\begin{enumerate}[nosep]
  \item \texttt{hd}, the head extractor;
  \item \texttt{tl}, the tail extractor;
  \item \texttt{resI}, the result $I$ arm;
  \item \texttt{resK}, the result $K$ arm;
  \item \texttt{Y}, the fixed-point combinator;
  \item \texttt{resS}, the result $S$ arm;
  \item \texttt{pair}, the pair constructor;
  \item \texttt{App}, the application constructor;
  \item \texttt{UQ}, the unquoter;
  \item \texttt{rb1}, the first rebuilder arm;
  \item \texttt{spApp}, the spine application handler;
  \item \texttt{rb}, the rebuilder;
  \item \texttt{stepI1}, the first step $I$ arm;
  \item \texttt{stepK2}, the second step $K$ arm;
  \item \texttt{stepI}, the step $I$ arm;
  \item \texttt{stepK1}, the first step $K$ arm;
  \item \texttt{stepK}, the step $K$ arm;
  \item \texttt{sp}, the spine walker;
  \item \texttt{stepS3}, the third step $S$ arm;
  \item \texttt{stepS2}, the second step $S$ arm;
  \item \texttt{stepS1}, the first step $S$ arm;
  \item \texttt{stepS}, the step $S$ arm;
  \item \texttt{step}, the step walker from the spine to the head leaf;
  \item \texttt{loop1}, the first loop arm; and
  \item \texttt{whnfF}, the whole interpreter.
\end{enumerate}
\end{multicols}

\begin{table}[htb!]
\caption{Every named part of the \texttt{whnfF} interpreter, ordered by 
size.  Above the rule, the closed $\{S,K,I\}$ term in full; below it, 
the same term named against the parts above by $\mathit{lift}$, which 
renames exact structural matches as a macro.}\label{tab:whnfF}
\centering\footnotesize
\begin{tabular}{@{}l r >{\raggedright\arraybackslash\ttfamily}p{4.35in}@{}}
\toprule
\textrm{part} & \textrm{atoms} & \textrm{expansion} \\
\midrule
\texttt{hd}     &   4 & S I (K K) \\
\texttt{tl}     &   5 & S I (K (K I)) \\
\texttt{resI}   &  11 & S (K K) (S (K K) (S (K (S I)) K)) \\
\texttt{resK}   &  13 & S (K K) (S (K (S (K K))) (S (K (S I)) K)) \\
\texttt{Y}      &  14 & S (K (S I I)) (S (S (K S) K) (K (S I I))) \\
\texttt{resS}   &  15 & S (K (S (K K))) (S (K (S (K K))) (S (K (S I)) K)) \\
\texttt{pair}   &  17 & S (S (K S) (S (K K) (S (K S) (S (K (S I)) K)))) (K K) \\
\texttt{App}    &  32 & S (K (S (K K))) (S (K (S (K K))) (S (K (S (K K))) (S (S (K S) (S (K K) (S (K S) (S (K (S I)) K)))) (K K)))) \\
\texttt{UQ}     &  43 & S (K (S I I)) (S (S (K S) K) (K (S I I))) (S (K (S (S (S (S I (K S)) (K K)) (K I)))) (S (K K) (S (S (K S) (S (K (S (K S))) (S (K K)))) K))) \\
\texttt{rb1}    &  44 & S (S (K S) (S (K K) (S (K S) K))) (K (S (K (S (K K))) (S (K (S (K K))) (S (K (S (K K))) (S (S (K S) (S (K K) (S (K S) (S (K (S I)) K)))) (K K)))))) \\
\texttt{spApp}  &  48 & S (S (K S) (S (K K) (S (K S) (S (K (S (K S))) (S (K K)))))) (K (S (K K) (S (K (S (S (K (S (K K))) (S (S (K S) (S (K K) (S (K S) (S (K (S I)) K)))) (K K))))) K))) \\
\texttt{rb}     &  74 & S (K (S I I)) (S (S (K S) K) (K (S I I))) (S (K (S (S (K S) (S (K (S I)) K)))) (S (K (S (K K))) (S (S (K S) (S (K K) (S (K S) K))) (K (S (K (S (K K))) (S (K (S (K K))) (S (K (S (K K))) (S (S (K S) (S (K K) (S (K S) (S (K (S I)) K)))) (K K))))))))) \\
\midrule
\texttt{stepI1} & 86 & S (K (S (K just))) rb \\
\texttt{stepK2} & 89 & S (K K) stepI1 \\
\texttt{stepI} & 92 & S hd (K stepI1) \\
\texttt{stepK1} & 99 & S (K (S hd)) (S (K K) stepK2) \\
\texttt{stepK} & 105 & S hd (K stepK1) \\
\texttt{sp} & 126 & Y (S (K (S (S (K S) (S (S (K S) (S (S (K S) (S (S (K S) K) (K resS))) (K resK))) (K resI))))) (S (K K) spApp)) \\
\texttt{stepS3} & 207 & S (K (S (K (S (K (S (K Just))))))) (S (K (S (K (S (K rb))))) (S (S (K S) (S (K K) (S (K S) (S (K (S (K App))) App)))) (K App))) \\
\texttt{stepS2} & 221 & S (K (S (K (S hd)))) (S (K (S (K K))) stepS3) \\
\texttt{stepS1} & 231 & S (K (S hd)) (S (K K) stepS2) \\
\texttt{stepS} & 237 & S hd (K stepS1) \\
\texttt{step} & 569 & S (S (S (S sp (K K)) (K stepS)) (K stepK)) (K stepI) \\
\texttt{loop1} & 588 & S (K (S (S (K S) (S (K K) (S step Just))))) K \\
\texttt{whnfF} & 618 & Y (S (K (S (S (K S) (S (K K) hd)))) (S (K K) loop1)) \\
\bottomrule
\end{tabular}
\end{table}

\noindent
Note that the expansions of \texttt{rb}, \texttt{UQ}, \texttt{sp}, and 
\texttt{whnfF} begin with the fourteen atoms of $Y$.  The fixpoint tie 
is a prefix of the emitted term, and which arms recurse can be read off 
the object itself.  This is perhaps the one piece of the surface's 
structure that survives bracket abstraction in a form a reader can 
observe for himself.

\begin{table}[htb]
\centering
\caption{The object terms of $SKI$ in the metacircular evaluator.}\label{tab:ski}
% \footnotesize
\begin{tabular}{@{}l l >{\raggedright\arraybackslash\ttfamily}p{3.35in} r@{}}
\toprule
\textrm{form} & \textrm{definition} & \textrm{expansion} & \textrm{atoms} \\
\midrule
\texttt{nil}            & $\lambda n.\lambda c.\,n$       & K & 1 \\
\texttt{cons}\;$h\;t$   & $\lambda n.\lambda c.\,c\,h\,t$ & S (K (S (K K))) (S (S (K S) (S (K K) (S (K S) (S (K (S I)) K)))) (K K)) & 22 \\
\texttt{nothing}        & $\lambda n.\lambda j.\,n$       & K & 1 \\
\texttt{just}\;$x$      & $\lambda n.\lambda j.\,j\,x$    & S (K K) (S (K (S I)) K) & 8 \\
\texttt{zero}           & $\lambda z.\lambda s.\,z$       & K & 1 \\
\texttt{suc}\;$n$       & $\lambda z.\lambda s.\,s\,n$    & S (K K) (S (K (S I)) K) & 8 \\
\bottomrule
\end{tabular}
\end{table}

\section{The language}\label{sec:language}

The \texttt{whnfF} function of \citet{Davis2027} introduced a number of
independent operational terms within its full expression, and these
suggest language features such as its loop, its comparison and dispatch,
and its atoms (\texttt{suc}, \texttt{zero}, \texttt{nil}, \texttt{just},
etc.).  These features were composed into larger machine terms, which
are similar to the supercombinators of \citet{Turner1979}.

SKIjack's features are written as a set of object terms in the 
metacircular evaluator; in essence, it is a macro language with a gently 
typed surface discipline.  They name and highlight the components of
\texttt{whnfF} and other recurrent supercombinators, augmented with
affordances to make the language expressive and composable.

\subsection{Design principles}

Since quotation of live terms is not supported, the language relies on 
the macro-like composition of its object terms rather than runtime 
evaluation.  A hard line between runtime quotation and compile-time
quotation is forced by Proposition~3.1 of \citet{Davis2027}, which follows from Theorem~3.2 of \citet{JayGivenWilson2011}:  ``No closed term $Q$ over $\{S,K,I\}$ satisfies $Q\,M \twoheadrightarrow \langle M\rangle$ for every closed $M$.''  Everything on the encoded side of 
that line is closed under the calculus, including running encoded terms, 
comparing them, and building cores.  Everything that would turn a live 
value back into data is unavailable and must be staged via the virtual
machine.  One design problem is to make that line visible and 
comfortable.  Another difference from traditional macro languages is 
that weak SKI reduction is call-by-name: arguments are never forced 
early, \texttt{K I Ω} is \texttt{I}, and no form needs thunking.

The basic components of SKIjack are thus:  a macro layer; Scott encoding;
recursion through $Y$; interpreters, the unquoter, and equality on
encoded terms like paths; the behavioral verification discipline; and
the existence of a runtime which can jet the compiled code for more
efficient execution.

SKIjack supplies one grammar with two lexicons (ASCII-only and Unicode), 
with a bijective token table checked at import, and the round-trip and 
lift laws as tests.  Given agent-driven development workflows, it is
pragmatic to maintain both depending on the aesthetic and practical 
considerations of the development environment.
Listing~\ref{lst:unicode} and Listing~\ref{lst:ascii} show the same 
simple example.

\begin{lstlisting}[label=lst:unicode,
                   caption={A simple example of SKIjack code under the
                   Unicode lexicon.}]
seg  ≡ Nat ∣ Two ∣ Three
path ≡ Nil ∣ Cons seg path
resolve ≔ ⦃/nat/two ↦ ⟨I⟩, /nat/three ↦ ⟨K⟩⦄
answer  ≔ wfN resolve ⊢ ⟨∵/nat/three⟩₁₀
\end{lstlisting}

\begin{lstlisting}[label=lst:ascii,
                   caption={A simple example of SKIjack code under the
                   ASCII lexicon.}]
seg  === Nat | Two | Three
path === Nil | Cons seg path
resolve := ns{/nat/two => <I>, /nat/three => <K>}
answer  := wfN resolve |- <?^/nat/three>@10
\end{lstlisting}

\subsection{Representation ABI}


\subsection{Kernel forms}

\subsection{Supercombinators vs.\ macros}


\subsection{The standard subject}

\subsection{The quotation line}




\begin{lstlisting}[float=htb,
                   label=lst:whnfF,
                   caption={The \texttt{whnfF} function defined in 
                   SKIjack.  Pleasingly, this program produces the 
                   identical $SKI$ artifact as Appendix B of
                   \citet{Davis2027} despite SKIjack not existing at the 
                   time of composition.}]
term === S | K | I | App term term
maybe === Nothing | Just term
whnfF := {
  step m = sp m nil stepS stepK stepI
  loop1 f m n2 = step m (Just m) (f n2)
  loop n m = n Nothing (loop1 loop m)
}
resS acc c0 c1 c2 = c0 acc
resK acc c0 c1 c2 = c1 acc
resI acc c0 c1 c2 = c2 acc
spApp f acc t u = f t (cons u acc)
sp m acc = m (resS acc) (resK acc) (resI acc) (spApp sp acc)
rb1 f h x xs = f (App h x) xs
rb h args = args h (rb1 rb h)
stepI1 x rest = Just (rb x rest)
stepI args = args Nothing stepI1
stepK2 x y rest = Just (rb x rest)
stepK1 x r = r Nothing (stepK2 x)
stepK args = args Nothing stepK1
stepS3 x y z rest = Just (rb (App (App x z) (App y z)) rest)
stepS2 x y r2 = r2 Nothing (stepS3 x y)
stepS1 x r = r Nothing (stepS2 x)
stepS args = args Nothing stepS1
\end{lstlisting}


SKIjack is not self-bootstrapping:  in any case, self-compilation would
require a parser which mapped arbitrary ASCII input to the internal 
representation of SKIjack terms.  While not impossible, the operational
overhead does not seem to merit the additional effort or reveal new
premises for $SKI$ as a tool.  Furthermore, quotation of live terms is 
not supported, which limits the ability to perform self-representation 
within the language itself (no runtime \texttt{eval}).  The language is
Turing-complete but certain functions are transferred to the VM,
such as scaffolding the initial evaluation and preserving the state and
memory arena.

\section{Execution model}\label{sec:execution-model}

\section{Interpreters as declarations}\label{sec:interpreters}

\status{The payoff section, and where the pair of papers is joined. The
declaration form and what it generates; then the one-site extension,
with \texttt{wf5Abs} and \texttt{wf5Omg} differing at a single arm and
compiling to 768 and 771 atoms, byte-identical to the artifacts; the T3
observability table reproduced. The sentence this section exists to
support: the first paper's move --- author a semantics at one site ---
is here a declaration form, and its cost over the hand construction is
zero atoms.}

\section{The namespace}\label{sec:namespace}

\status{Typed paths as a declared type rather than arbitrary terms; the
resolver as a core parameter; the namespace literal and what it compiles
to; blocking, the resume loop, and why re-running from scratch is what
makes it sound under an append-only namespace; the value-lookup example.
State the two honest wrinkles in the text and not a footnote: the
monotonicity hypothesis, and that fuel is not neutral across a resume.}

\section{Boundary soundness}\label{sec:boundary}

\status{Recommended, and the paper's own theorem rather than a borrowed
one. Define the fragment; state that an accepted program never places a
function where data is required, so the operations the calculus cannot
perform are refused before they can be reached; prove by induction on
the checker. Add erasure as a corollary --- checking can reject but
never rewrites, so the emitted term is the same whether it ran or not
--- which is what keeps the laws of \S\ref{sec:impl} intact. If this
section does not survive, say plainly in \S\ref{sec:intro} that the
enforcement is tested and not proved.}

\section{Implementation and evaluation}\label{sec:impl}

\status{The expander as a fold, its passes, and the dictionary. The four
laws as executable tests. The byte-identity table --- \texttt{whnfF}
618, \texttt{wf5Abs} 768, \texttt{wf5Omg} 771, \texttt{wfQ} 950,
\texttt{wfN} 1{,}066, \texttt{EQ5} 240, with T0 at 340 --- and, in the
same table, the distinction that answers the circularity objection
before it is raised: the generated forms were designed to reproduce the
artifact, whereas the long version, written out as ordinary equations
with generation switched off, reaches the same term without that help.
The corpus and the test count. One sentence stating that no performance
claim is made here and why: the reference host is an oracle, not an
implementation, and the VM is discussed only briefly.}

\section{The virtual machine}\label{sec:runtime}

\section{Limitations, and what is not claimed}\label{sec:limits}

\status{Self-hosting, declined, with the reason stated precisely: the
expander's input is a syntax tree, so it is separable from the parser,
and what is forgone is not a demonstration but the Futamura projections
inside the calculus. The type discipline stops at Stage A. There is no
standard subject yet, so a level-1 library name is inlined rather than
shared. And the precision this section
must carry: the language has evaluation of encoded terms, including
terms built at run time, and lacks reification of live values --- the
missing operation is \emph{reify}, not \emph{eval}. The paragraph
currently at the end of \S\ref{sec:language} belongs here.}

\bibliographystyle{plainnat}
\bibliography{mss}

\appendix

\section{The artifact}\label{app:artifact}

\status{The parts below the rule of Table~\ref{tab:whnfF}, in full, and
\texttt{whnfF} itself at 618 atoms and 2{,}109 characters.  The first
paper prints the same term in its Appendix B; if space is short, cite it
there and print only the parts this paper names, since the claim being
supported is byte-identity and a reader who wants to check it needs the
term from one of the two.}


\end{document}
